\documentclass[10pt, twocolumn]{article}

\usepackage[utf8]{inputenc}
\usepackage{amsmath, amssymb}
\usepackage{geometry}
\geometry{letterpaper, margin=0.65in}
\usepackage{graphicx}
\usepackage{booktabs}
\usepackage{enumitem}
\usepackage{url}
\usepackage[hidelinks]{hyperref}
\usepackage[compact]{titlesec}
\titlespacing{\section}{0pt}{1.6ex plus .2ex}{1.0ex plus .2ex}
\titlespacing{\subsection}{0pt}{1.3ex plus .2ex}{0.7ex plus .2ex}

\newcommand{\code}[1]{\texttt{\small #1}}
\setlength{\abovedisplayskip}{4pt}
\setlength{\belowdisplayskip}{4pt}
\setlength{\textfloatsep}{7pt}
\setlength{\intextsep}{7pt}

\title{\textbf{Reward, Reset, and Evaluation Design for\\Sim-to-Real Reinforcement Learning:\\A Case Study in Quadrotor Gate Racing}}
\author{Zhuohao Ni\thanks{These authors contributed equally to this work.}
\and Bruce (Yuqian) Zhang\footnotemark[1]}
\date{}

\begin{document}

\maketitle

\begin{abstract}
\small\noindent
We report a two-phase reinforcement learning project in which one quadrotor
racing task---a seven-waypoint course containing a self-intersecting
double-gate structure---was first solved as a pure simulation problem and then
re-solved as a transfer problem targeting a physical vehicle tracked by motion
capture. The phases share a track, a simulator and a learning algorithm, and
almost nothing else: the reward, the observation space, the reset distribution,
the randomization strategy and the checkpoint-selection rule were each
re-derived, and in four of five cases inverted, once the simulator stopped
being the objective and became a hypothesis. Dense shaping with hand-placed
virtual targets reached $97.3\%$ three-lap success at $15.68$\,s mean over
$5000$ randomized environments; a sparse gate-centric reward with an
observation space constrained by the deployment interface reached $98.0\%$ at
$18.34$\,s under a re-identified dynamics nominal and completed the physical
three-lap evaluation in $19.004$\,s, continuing to five correctly ordered laps.
Three design objects usually treated as implementation details carried most of
the transfer risk for us: the reward predicate's relationship to the task's
progress predicate, the reset distribution as an explicit choice of where the
value function is fit, and the evaluation harness---which failed in both
directions, once by being more diverse than deployment and once by being
calibrated so closely to measured hardware that it ranked first a policy that
managed a single lap in the arena.
\end{abstract}

\section{Introduction}

Reinforcement learning for agile flight is usually presented as a transfer
problem: train in simulation, randomize the physics, deploy
\cite{loquercio2020,kaufmann2023,song2023}. Framed that way the interesting
variable is the randomization distribution and the reported result is a success
rate. The framing hides the fact that a policy trained against a simulator
whose fidelity is not in question and a policy trained against a simulator
standing in for hardware solve different optimization problems, even when the
MDP is written identically.

We solved both. The task is fixed: three laps of a seven-waypoint gate course,
in order, without crashing. In the \emph{simulation phase} correctness was
defined by the simulator, so the rational strategy was to shape as aggressively
as training stability allowed. In the \emph{sim-to-real phase} that simulator
became a model with unknown parameters, an observation pipeline with a biased
physical counterpart, and evaluation numbers that had to predict an outcome
measurable only a handful of times. Four decisions inverted:

\begin{enumerate}[nosep, leftmargin=1.4em]
\item \textbf{Dense $\rightarrow$ sparse reward.} Shaping that specifies
\emph{how} to fly encodes a trajectory manifold optimal only under the training
dynamics.
\item \textbf{Exploit tolerance $\rightarrow$ exploit detection.} In simulation
we removed a termination rule punishing reverse gate passes, because tolerating
the mistake trained more smoothly; in transfer we added three detectors.
\item \textbf{Randomization as insurance $\rightarrow$ randomization as
identified uncertainty.} Widening the distribution around an unidentified
nominal buys robustness in directions that do not matter.
\item \textbf{Training return $\rightarrow$ deployment metric for checkpoint
selection.} The standard \code{best\_model.pt} rule---argmax of mean training
return---selected a checkpoint whose measured three-lap success rate was
$84.7\%$.
\end{enumerate}

The fifth change was not an inversion but a constraint: the observation space
stopped being a free design variable and became an interface contract with the
deployment controller.

Our contributions are empirical rather than algorithmic: we use standard PPO
\cite{schulman2017} with GAE \cite{schulman2016} throughout, and the learning
algorithm is not where the difficulty lay. We offer a measured account of (i) a
reward-hacking taxonomy induced by a self-intersecting task graph and a general
fix, (ii) the reset distribution as a first-class design object including its
interaction with actuator state, (iii) a sensitivity analysis that reallocated
our randomization budget from dynamics to observations, and (iv) two opposite
failure modes of evaluation harness design.

\section{Problem Setup}

\textbf{Task and control stack.} The vehicle is a Crazyflie-class quadrotor
simulated in Isaac Lab \cite{mittal2023}. Rigid-body physics runs at $500$\,Hz;
the policy acts at $50$\,Hz (decimation $10$) and emits
$a_t = [\,\bar{T}, \omega_x^{\text{des}}, \omega_y^{\text{des}},
\omega_z^{\text{des}}\,] \in [-1,1]^4$, a normalized collective thrust and
three body-rate commands, which an inner $500$\,Hz PID rate controller converts
to a desired wrench tracked by a first-order motor model. Episodes last $30$\,s
or terminate on crash, capped at three laps. This action interface is the one
exposed by the physical stack, and it fixes what \emph{can} be randomized:
everything between action and wrench.

\textbf{Track geometry.} Six waypoints are ordinary; the interesting structure
is central. Waypoints 2 and 3 are gate planes $1.25$\,m apart at $z = 0.75$\,m
with identical heading, and waypoint 6 is \emph{the same plane as waypoint 3,
entered from the opposite side}. The course self-intersects: the state
(position, velocity) does not determine which waypoint the agent is pursuing,
and one physical gate crossing can be a legitimate pass, an illegitimate
reverse pass, or a pass toward a different waypoint index, depending on
history. The remaining waypoints sit at $(\pm 1.5, \pm 3.5, 2.0)$ and
$(2.0, \pm 3.5, 0.75)$, so each lap couples two wide high-altitude turns to a
tight low-altitude transit of the central structure.

Nearly every difficulty in this project traces to that fact. It makes gate-pass
detection a stateful predicate rather than a geometric test; it creates a
family of reward exploits (Section~\ref{sec:hacking}); and because the drone
must gain and lose $1.25$\,m of altitude in the shortest segment, it is where
dense shaping and dynamics mismatch bite hardest. A pass at waypoint $i$ is detected in the
gate's local frame---the agent's $x$ crossing from positive to negative while
$(y,z)$ lies inside an opening of half-side $s$.

\textbf{Evaluation protocols.} We report four. \emph{Self-evaluation}
(simulation phase): batch rollouts over thousands of parallel environments,
each drawing independent perturbations of thrust-to-weight, drag and rate-loop
gains. \emph{External evaluation} (simulation phase): an independently
specified protocol with fixed perturbation bounds and a reference time of
$16.06$\,s applied to a submitted policy; we report our own measurements
against its bounds rather than its outputs. \emph{Batch evaluation}
(sim-to-real phase): reset-diverse rollouts under the re-identified dynamics
nominal, our primary robustness screen. \emph{Physical evaluation}: timed
three-lap runs on the real vehicle, with lap timing recovered post hoc from
flight logs under a strict gate-ordering rule. \textbf{These protocols are not
mutually comparable, and we do not compare them.} Simulation-phase and
sim-to-real-phase times differ by roughly $2.7$\,s not because the policy got
worse but because the phases run under different dynamics nominals ($3.15$
versus the value identified from flight logs, Section~\ref{sec:twr}). A policy
with less thrust authority is slower; that is physics, not regression.

\section{Learning Algorithm and Infrastructure}

We use PPO \cite{schulman2017} with the clipped surrogate objective, an entropy
bonus, and GAE \cite{schulman2016} ($\gamma = 0.99$, $\lambda = 0.95$) across
$8192$--$16384$ parallel environments in the style of massively parallel
on-policy learning \cite{rudin2022}. Early prototypes collapsed around $2500$
iterations, which a KL-adaptive learning rate ($\bar{d}_{\mathrm{KL}} = 0.01$)
removed; it did not remove the reward-driven collapses of
Section~\ref{sec:sim-reward}, which were a shaping problem rather than a
step-size problem.

Table~\ref{tab:ppo} lists the configuration by phase. Two entries are
consequences of decisions made later. The rollout length increase from $24$ to
$64$ steps is forced by the reward becoming sparse: at $50$\,Hz a $24$-step
rollout spans $0.48$\,s, which frequently contains no gate-pass event, so most
advantage estimates carry no task signal. And \code{empirical\_normalization}
is disabled in transfer not because it hurts but because the running
observation normalizer is a \emph{learned component of the policy} the
deployment controller never receives; hand-chosen fixed constants failed to
recover the same performance, which we read as evidence the running statistics
were doing real work---work we then had to forgo. The deeper critic follows
Pasumarti et al.\ \cite{pasumarti2025}---value estimation under sparse racing
objectives benefits from more capacity than the policy needs---while actor
width is fixed by the deployment controller.

\begin{table}[t]
\centering
\footnotesize
\caption{PPO and network configuration by phase.}
\label{tab:ppo}
\begin{tabular}{@{}l r r@{}}
\toprule
& Simulation & Sim-to-real \\
\midrule
Actor hidden dims    & $[256,256]$ & $[512,512,256,128]$ \\
Critic hidden dims   & $[512,256,$ & $[512,512,256,$ \\
                     & $128,128]$  & $256,128,128]$ \\
Steps per env        & $24$ & $64$ \\
Mini-batches         & $4$ & $8$ \\
Learning rate        & $5\!\times\!10^{-4}$ & $1\!\times\!10^{-4}$ \\
Entropy coefficient  & $0.005$ & $0.01$ \\
Empirical obs. norm. & enabled & \textbf{disabled} \\
\bottomrule
\end{tabular}
\end{table}

\section{Simulation Phase}

\textbf{Observation space.} The policy receives $31$ dimensions, spatial
quantities in the body frame: body linear and angular velocity ($6$); gravity
projected into the body frame ($3$); current, next and next-next gate positions
($9$); entry normals for the current and next gates ($6$); normalized waypoint
index ($1$); $\sin$ and $\cos$ of yaw error ($2$); previous action ($4$). The
gravity vector replaces the unit quaternion, preserving all tilt information
the task needs without the double-cover sign ambiguity. The waypoint index is a
deliberate concession to self-intersection: since position and velocity do not
identify the active waypoint, the index restores the Markov property the
geometry destroys, and without it the policy cannot in principle distinguish
waypoint 3 from waypoint 6.

\subsection{Reward and the collapse modes it produced}
\label{sec:sim-reward}

Table~\ref{tab:reward} gives the converged configuration alongside its
sim-to-real successor. It is the result of a search whose failures are more
informative than its endpoint.

The instructive failure was a heading-alignment term ($\cos$ of yaw error)
added alongside a velocity reward at scale $5.0$, whose episodic contribution
reached ${\sim}750$ against the gate term's ${\sim}200$. The policy became
speed-dominated: it accelerated toward each gate at maximum rate and could not
decelerate for the tight transition out of the central structure, gate reward
collapsing from ${\sim}200$ to ${\sim}40$ while velocity reward rose to
${\sim}150$. The relevant comparison is not between reward \emph{scales} but
between \emph{integrated episodic contributions}, and a dense term integrated
over $30$\,s can silently dominate a sparse term worth $200$ per event.

\begin{table}[t]
\centering
\footnotesize
\caption{Reward terms by phase. Progress, velocity blending, orientation
penalty and the virtual-target machinery are all removed in transfer; the gate
bonus is unchanged and does all the work.}
\label{tab:reward}
\begin{tabular}{@{}l r r@{}}
\toprule
Term & Sim & Sim-to-real \\
\midrule
Gate pass (indicator)          & $200.0$ & $200.0$ \\
Progress ($\Delta$dist, clamped) & $50.0$  & off \\
Racing-line velocity (blended) & $8.0$   & --- \\
Orientation ($|\phi|+|\theta|$ over $0.5$)  & $-1.0$  & --- \\
Lap-incomplete (per step)      & ---     & $-0.05$ \\
Rate reg.\ roll/pitch ($\omega_{x,y}^2$) & --- & $-1.0$ \\
Rate reg.\ yaw ($\omega_z^2$)  & ---     & $-0.5$ \\
Action smoothness              & $-0.2$  & $-0.1$ \\
Contact (per step) / death     & $-1.0/-10$ & yes \\
\bottomrule
\end{tabular}
\end{table}

\subsection{Shaping as implicit path planning}

Rather than specify a geometric racing line, we shaped the velocity reward to
induce one. With $\hat{d}_{\text{cur}}, \hat{d}_{\text{next}}$ the unit vectors
to the current target and the following gate,
\begin{equation}
r_{\text{vel}} \propto
\alpha \, \langle v, \hat{d}_{\text{cur}} \rangle +
(1-\alpha) \, \langle v, \hat{d}_{\text{next}} \rangle ,
\end{equation}
with $\alpha = 5/8$ on ordinary gates. Because part of the reward points at the
\emph{next} gate before the current one is reached, the optimal trajectory cuts
toward the inside of each corner---smooth curving racing lines with no path
parameterization, no spline and no per-corner tuning.

The blend is a genuine optimization knob with a stability boundary. Reducing
$\alpha$ to $4/8$---an apparently mild change---caused late-training collapse
at ${\sim}3000$ iterations, three-lap success dropping from ${\sim}90\%$ to
zero and mean time spiking to ${\sim}20$\,s before recovering to a degraded
plateau. The next-gate component pulls the vehicle far enough off the current
gate's approach at high speed to cause misses, and once misses are common the
sparse gate reward degrades faster than the dense velocity term---a
self-reinforcing dynamic. We report this because $\alpha$ looks like a
hyperparameter and behaves like a phase boundary.

\textbf{Virtual targets.} Because the gate-pass predicate reads the pose in the
gate frame, it is independent of whatever position the reward aims at---so
shaping can be steered without altering the task definition. We exploited this
heavily: the central structure was designed to force a vertical loop, a path
around its \emph{side} proved faster, and we induced that path with a
two-phase virtual target placed off the gate centers entirely. The technique is
powerful and does not survive contact with hardware, since it encodes a
specific solution to a specific geometry under specific dynamics.

\subsection{Reset distribution}

Resets draw the target waypoint uniformly over all seven, so every segment
trains equally, and place the vehicle behind it at a distance ramping
$[1.0,2.0] \rightarrow [1.5,4.0]$\,m over $800$ iterations, with lateral,
vertical and yaw noise; $40\%$ of episodes instead start mid-segment, uniform
at $20$--$80\%$ of the way between consecutive gates, ${\sim}20\%$ start on the
ground, and half of airborne spawns receive $0.5$--$3.0$\,m/s of forward
velocity.

One entry deserves emphasis: raising the mid-segment rate from $25\%$ to $40\%$
moved three-lap success from $86.7\%$ to $96.3\%$. This was the single largest
measured improvement of the simulation phase, larger than any reward change,
and it was obtained by altering where episodes begin rather than what they
reward. The mechanism is straightforward once stated: on-policy methods fit the
value function only on the state distribution induced by $\rho_0$ and the
current policy, so spawning behind gates trains approaches well and
gate-to-gate \emph{transitions} poorly, since transition states are visited
only by trajectories that already succeeded at everything preceding them.

\subsection{Randomization, paired evaluation, results}

Thrust-to-weight, drag coefficients and rate-loop PID gains were resampled per
environment at every reset, following standard dynamics randomization practice
\cite{tobin2017,peng2018}. The final model was a baseline trained $5000$
iterations at $16384$ environments, then fine-tuned $2000$ more under widened
bounds (thrust-to-weight $\pm8\%$, drag $0.3$--$2.5\times$, PID $\pm25\%$ on
$k_p,k_i$ and $\pm40\%$ on $k_d$) with entropy reduced from $0.005$ to $0.001$.

The most instructive comparison was a \emph{paired} evaluation: $1000$
environments with identical perturbation draws presented to both models. Rank
correlation between per-environment times was essentially zero (Pearson
$r = -0.016$, Spearman $\rho = 0.048$). Fine-tuning under wider bounds did not
translate the performance distribution; it \emph{rotated} it. Bucketed by
baseline time, the fine-tuned model was $0.87$\,s \emph{slower} in the
baseline's best bucket ($<15.6$\,s) and $1.32$, $0.99$ and $1.38$\,s faster in
the borderline, at-risk and worst buckets, converting $168$
environments---$78.5\%$ of the baseline's failures---into passes while
introducing $137$ new ones, a net gain of $31$. A headline success-rate
improvement of $2.4$ percentage points therefore concealed $305$ environments
whose outcome changed, in both directions. Aggregate metrics on independently
randomized environments cannot distinguish a uniform improvement from a
reshuffle, and a reshuffle means the new model has blind spots about which the
old model's evaluation history says nothing.

Under self-evaluation across $5000$ independently randomized environments the
fine-tuned model achieved $97.3\%$ three-lap success at $15.68$\,s (median
$15.58$, $\sigma = 0.50$, best $14.82$, worst $21.26$\,s), against the
baseline's $94.9\%$ and $15.78$\,s ($\sigma = 0.54$). Measured against the
external protocol's $16.06$\,s reference, $84.0\%$ of environments finished
inside it, against $82.0\%$ for the baseline.

\section{Sim-to-Real Phase}

Transfer changed what was adjustable before it changed what was optimal.

\subsection{The observation space becomes a contract}

The deployment controller computes its own observation vector from
motion-capture data, so any quantity the policy consumes must be reproducible
by that code. This eliminated most of the simulation-phase observation design:
engineered features---yaw error, waypoint index, gate normals---gave way to raw
geometry,
\begin{equation*}
\underbrace{v_b}_{3} \;\big|\;
\underbrace{\mathrm{vec}(R_{b \rightarrow w})}_{9} \;\big|\;
\underbrace{C^{\text{cur}}_b}_{12} \;\big|\;
\underbrace{C^{\text{next}}_b}_{12} \;\big|\;
\underbrace{a_{t-1}}_{4}
\end{equation*}
for $40$ dimensions, where $C_b$ is a gate's four corners in the body frame,
computed identically on both sides as $C_w = S R_g^\top + p_g$ then
$C_b = (C_w - p_d) R_b$. The full rotation matrix replaces the gravity vector
because, having given up the yaw-error feature, the policy needs complete
attitude. Twelve numbers per gate redundantly encode a four-degree-of-freedom
object, and we adopted them only because the deployment controller already
computed them: \emph{a feature that cannot be reproduced by the deployment
stack is not available at any price}, and feature engineering that helps in
simulation becomes a liability once the pipeline must exist twice and agree.
Including the previous action follows Kaufmann et al.\ \cite{kaufmann2023}:
under actuator lag it carries information about actuator state that the
kinematic observation does not.

\subsection{Sparse reward redesign}

Progress, velocity blending, orientation penalty and virtual targets were all
removed (Table~\ref{tab:reward}). The argument, following the sparse-objective
results of \cite{pasumarti2025} and the observation that shaping alters the
optimal policy unless potential-based \cite{ng1999}, is that dense terms
specify \emph{how} to fly and thereby encode a trajectory manifold optimal only
under the training dynamics; when those dynamics shift, nothing has given the
policy reason to have learned anything off that manifold. Two results
generalize.

\textbf{Sparse bonuses must beat integrated, not instantaneous, costs.} Our
first sparse configuration set the gate bonus to $10.0$ and collapsed at
${\sim}200$ iterations: the policy learned to hover. With a constant per-step
penalty and a per-step contact penalty always active, and a gate taking on the
order of a second to reach, the expected integrated cost of \emph{attempting} a
gate exceeded the bonus for succeeding; raising the bonus to $200.0$ fixed it.
A sparse bonus must exceed the integrated negative cost accumulated over the
expected time to earn it---for a $30$\,s episode at $50$\,Hz, a factor of
hundreds away from the instantaneous comparison one is tempted to make.

\textbf{Command regularization is a calibrated speed dial.} We ablated the
roll/pitch and yaw rate penalties from $(-1.0,-0.5)$ to $(-2.0,-0.05)$, since
flight logs show yaw rate staying low while roll and pitch drive the aggressive
transients, so the default split under-penalized the axis producing jerk. The
rebalanced policy was the smoothest we trained---lowest mean tilt and body rate
of any policy we flew---and reliably slower by ${\sim}0.6$\,s per lap. This is
an exchange rate, not a failure: the term buys attitude margin at a measurable
price in time, and it is the correct knob when margin binds.

\subsection{Specification gaming on a self-intersecting task graph}
\label{sec:hacking}

Because waypoints 2, 3 and 6 share the central structure and waypoint 6 reuses
waypoint 3's plane in reverse, a naive pass predicate admits several ways to
earn credit without flying the intended course. We met three.
\emph{Reverse crossing of the current target}: the vehicle crosses its target
plane the wrong way, detected by the mirrored predicate and treated as
termination, with a five-step cooldown after a legitimate pass since frame
recomputation at the moment of passing otherwise produces false positives.
\emph{Plane-straddling at the loop gate}: the vehicle crosses waypoint 3's
plane by a few centimetres, so $x_{t-1} > 0$ and $x_t \leq 0$, then immediately
reverses out, collecting the bonus without traversing the structure. The fix
conditions the pass on having genuinely been in the approach zone---we track
$\max_{\tau \leq t} x_\tau$ since the target last changed and require it to
exceed a threshold---applied \emph{only} at waypoint 3, because applying it
globally caused visibly valid passes elsewhere to go uncounted, corrupting our
own evaluation. \emph{Backtracking through the previous gate}: after passing
waypoint 2 the vehicle reverses back through its plane to re-enter the $+x$
approach region, then proceeds to waypoint 3---a cheap way to restart the
approach, available precisely because waypoints 2 and 3 share a structure and a
heading, and invisible to checks that monitor only the \emph{current} target.
The fix is a generic previous-gate detector on waypoint $(i-1) \bmod 7$, which
needs care at the boundaries---it must not fire before any gate is passed, nor
on the step a gate is passed, when the previous-gate index has just changed and
the buffered comparison is meaningless, nor outside the opening.

\textbf{Splitting the reward predicate from the progress predicate.} Those
fixes are patches; the structural change that made them coherent was to stop
using one predicate for two purposes. Earlier configurations used a single
half-side $s$ for everything, so shrinking $s$ to encourage center-line flight
also shrank the region in which the agent's \emph{target advanced}, producing
episodes where the vehicle physically flew through a gate and the task did not
acknowledge it. That is not conservatism; it is a corrupted task definition. We
now maintain two predicates over the same crossing event: a \textbf{progress
predicate} at the physical half-side ($0.5$\,m), which advances the target
index, increments the lap counter and drives termination; and a \textbf{reward
predicate} at a virtual inner half-side ($s/2 = 0.35$\,m), which alone controls
the gate bonus and insertion into the replay buffer of
Section~\ref{sec:reset}. The agent is paid for flying through the middle while
the task counts what physically happened. We regard this as the most portable
idea in this report: whenever a shaping term and a task-progress condition are
computed from the same geometric test, tightening the shaping silently
redefines the task. For the same reason we can zero dense progress at selected
waypoints, since Euclidean distance to a gate center shrinks as readily from
the \emph{wrong} side as the right one.

\textbf{The same exploit, two correct answers.} In the simulation phase we
experimented with strict termination on reverse and wrong-side passes and
\emph{deliberately removed it}: letting the agent make the mistake without
instantly ending the episode produced noticeably smoother training. In transfer
we added three detectors. Both were right for their regime, and the difference
is what the exploit costs. Where the simulator defines correctness, an exploit
costs a suboptimal trajectory while the termination suppressing it costs
gradient signal and stability---frequently a bad trade. Where the simulator
stands in for hardware, an exploit costs a policy that has learned a behavior
the real geometry will not support, discovered at deployment. The same
observation licenses opposite interventions depending on whether simulator
fidelity is in question, which we think is under-appreciated in discussions of
specification gaming \cite{amodei2016} that treat exploits as unconditionally
undesirable.

\subsection{The reset distribution as a design object}
\label{sec:reset}

Having found that reset placement outweighed reward tuning, we treated $\rho_0$
explicitly as a mixture of four components.

\emph{Gate-pass replay ($30\%$).} Following the practice of resetting to states
observed during successful execution \cite{kaufmann2023,song2023}, we maintain
a per-waypoint ring buffer keyed by the \emph{next} target gate: when the
reward predicate fires, the full root state and previous action are written to
the buffer for the newly acquired target, and at reset a state is drawn and
perturbed by $\pm0.15$\,m horizontally, $\pm0.10$\,m vertically,
$\pm0.20$\,m/s linear and $\pm0.50$\,rad/s angular velocity, $\pm0.20$\,rad
yaw. The buffer holds only states the policy actually reached, so the
distribution tracks the policy as it improves; we ramp its fraction up from
zero, since an unrepresentative buffer early in training supplies states the
policy cannot yet handle.

\emph{Actuator-state reconstruction.} A detail we consider necessary rather
than optional: replaying a rigid-body state into a system with actuator
dynamics does not replay the state. The rotor speeds at the recorded instant
were not at rest, and initializing them to zero injects a spin-up transient
into precisely the states we selected for being important. We therefore
recompute the desired wrench from the replayed previous action, invert the mixer
for motor-speed targets, and initialize both commanded and actual motor speeds
to it, so the first step after reset is dynamically continuous with the
trajectory the state came from.

\emph{Ground starts ($20\%$).} Physical evaluation begins at rest on the floor,
so a fraction of episodes start there, a few centimetres up and up to $3$\,m
behind the first gate, with a narrowed variant matching the launch pose
observed in the arena. \emph{Geometric fallback.} Remaining episodes
interpolate between consecutive gates with $t \sim \mathrm{Beta}(0.5,0.5)$,
concentrating mass near gates and thinning it mid-segment.

\emph{A cautionary ablation.} The fallback spawns at rest, unlike race flight,
so we added an option to give it velocity toward the target. At a cap of
$1.2$\,m/s the policy became measurably faster but lost success rate and grew a
worse time tail; at $0.6$\,m/s success recovered ($98.3\%$) but the speed gain
vanished entirely (mean $19.46$\,s against $18.34$\,s for the base). We had
framed this as distribution matching; the result says otherwise. Rather than
making training states more realistic, the higher reset velocity trained a more
aggressive state distribution---buying speed by shifting the policy, not by
correcting a mismatch.

\subsection{Randomization as identified uncertainty}
\label{sec:dr}

The environment randomizes the transition model---thrust-to-weight, drag,
rate-loop $k_p,k_i,k_d$, mass, motor time constant, action latency---and the
observation map: per-channel Gaussian noise ($\sigma = 0.05$\,m/s on body
velocity, $0.01$ on rotation entries, $0.02$\,m on gate corners), observation
latency, and systematic biases on body-frame velocity, gate-corner position and
yaw. The distinction between the two groups mattered more than any individual
range.

\textbf{Sensitivity analysis as budget allocation.} Before committing budget we
ran a deterministic evaluation-only sweep over thirteen perturbation scenarios,
$768$ ground-start episodes each, on a single checkpoint, asking whether this
policy was more likely to fail from moderate control mismatch or from
systematic observation mismatch. Every transition-model perturbation left
three-lap success at or above $95.6\%$, including $10\%$ thrust error, $40$\,ms
command latency and $\pm15\%$ rate-gain error, and zero-mean observation noise
at the trained magnitude was likewise harmless ($96.5\%$). But a constant
$0.20$\,m/s bias on observed body-frame velocity reduced success to $2.9\%$ at
a mean of $28.15$\,s, and a $5$\,cm systematic offset on gate corners cost
$4.2$ points and $2.1$\,s. Takeoff and first-gate success were $100\%$ in all
thirteen scenarios; the separating signal is entirely in later trajectory
execution.

The interpretation we favor is that a closed loop rejects transition-model
error as a matter of course---that is what feedback is for---whereas a
systematic observation bias is an unmodeled, unobservable state offset the
policy cannot detect and has no mechanism to correct. Randomizing dynamics
teaches robustness to a disturbance the controller was already going to reject;
randomizing observation \emph{bias} teaches robustness to something it cannot.
This is a sim-side result for one checkpoint, and several control scenarios
scored slightly \emph{above} nominal, which is run-to-run variation; what
survives is the ordering.

\textbf{Re-centering the transition model.}
\label{sec:twr}
Randomization defends against parameters near the center of its distribution.
It cannot defend against a center that is wrong. Flight logs indicated an
effective mass near $64$\,g against the simulator's nominal $30$\,g, while the
deployed controller maps full policy thrust to approximately $1.174$\,N, so
\begin{equation}
\mathrm{TWR}_{\text{real}} \approx
\frac{1.174\ \mathrm{N}}{0.064\ \mathrm{kg} \times 9.81\ \mathrm{m/s^2}}
\approx 1.87 ,
\end{equation}
against a simulator nominal of $3.15$. Our widest randomization to that point
was $\pm8\%$, spanning $[2.90,3.40]$. \textbf{The physical operating point was
not in a low-density tail of the training distribution; it was outside its
support entirely.} No width around $3.15$ would have reached it, and no reward,
architecture or reset change could have compensated: the policy was learning to
exploit control authority the vehicle did not have.

Re-centering on $1.87$ with a narrow $\pm6\%$ band produced $98.0\%$ success at
$18.34$\,s ($\sigma = 0.48$) over $768$ environments. Notably an
\emph{untrained-for} sanity check---the old checkpoint under the new
nominal---already completed three laps at $21.16$\,s, which gave us confidence
before committing to a retrain. The cost is specialization: the re-centered
policy performs poorly at the original nominal, terminating around seven gates
on representative rollouts. This is the robustness--specialization trade
documented by Ferede et al.\ \cite{ferede2025}---less randomization transfers
faster when the nominal is correct and fails outright when it is not---and we
accepted it because the deployment thrust scale is fixed and known, at the
price of a policy for \emph{this vehicle} rather than for quadrotors. Two
general statements follow: randomization widens support but does not move the
center, and the two are not interchangeable; and the cheapest diagnostic for a
suspected support failure is not a wider sweep but a single rollout at the
hypothesized nominal. Tightening the remaining ranges once the nominal was
corrected (thrust-to-weight $\pm4\%$, drag $0.8$--$1.2\times$, PID
$\pm8\%/\pm15\%$) did not buy speed---across four candidates the best reached
$98.3\%$ at $19.46$\,s---which suggests the residual gap for this policy family
is geometric, racing line and clearance, rather than a width effect.

\section{Evaluation and Policy Selection}
\label{sec:eval}

Three evaluation-design failures cost us more than any modeling error.

\textbf{Training return is not the deployment objective.} The framework saves
\code{best\_model.pt} at the argmax of mean episodic training return. Nothing
about that quantity optimizes what deployment requires: three-lap success rate,
low variance, a bounded worst case. In one run the saved best model was written
at ${\sim}$iteration $513$ at mean return $5746.86$; its batch evaluation was
$84.7\%$ success, $19.09$\,s mean, $\sigma = 1.26$\,s, $28.44$\,s worst
case---a checkpoint that wins on return by flying fast and occasionally failing
badly, exactly the profile training return cannot distinguish from a good one.
We replaced the rule with an explicit procedure: assemble a candidate pool
spanning the saved best model, the highest-return checkpoints and late-training
anchors, evaluate each under the deployment-matched protocol, and accept only if
SR $\geq 98.5\%$, mean $\leq 18.65$\,s, $\sigma \leq 0.5$\,s and no
pathological tail.
\textbf{A faster best case is never sufficient evidence}---a rule adopted only
after a candidate that won on best-case time failed in the arena.

\textbf{Failure one: an evaluation distribution deployment never samples.}
Because our batch evaluation runs with training-mode hooks active, it inherits
the \emph{training} reset mixture: at a ground-start fraction of $0.2$ a
nominally $1000$-environment evaluation contained roughly $210$ ground starts,
the other $790$ beginning mid-course, at altitude, from the buffered state of a
policy that had already succeeded---while every physical run begins on the floor
at rest. The fix is a deployment-matched profile forcing that fraction to
unity.

\textbf{Failure two: an evaluation calibrated past the point of usefulness.}
Having identified real parameters, the natural next step was a profile that
used them: identified thrust-to-weight, a fixed three-step action delay
matching the measured $60$--$70$\,ms command lag, ground starts only, the
narrow launch pose, and observation noise reduced from the training $0.05$\,m/s
to the measured stationary $0.005$\,m/s. Since a physical session is one
vehicle on one day, this seemed strictly more faithful. Backtested against
three policies whose real-world ranking we knew (Table~\ref{tab:backtest}), it
ranked \emph{first} the policy that in the arena completed a single correctly
ordered lap with negative clearance at a gate, and gave all three $100\%$
success with $\sigma < 0.1$\,s; a latency-stress variant exposed fragility but
over-penalized a policy that had completed four ordered laps. The older, less
faithful, reset-diverse screen correctly kept the base policy first.

\begin{table}[t]
\centering
\footnotesize
\caption{Backtest of evaluation profiles against known physical outcomes.
Real ranking: base $>$ speed-v4 $>$ speed-v2.}
\label{tab:backtest}
\begin{tabular}{@{}l r r r@{}}
\toprule
Policy & Calib.\ & Latency & Reset-diverse \\
       & mean (s) & stress SR & SR / mean (s) \\
\midrule
base     & $18.24$ & $100.0\%$ & $98.0\%$ / $18.34$ \\
speed-v2 & $\mathbf{17.96}$ & $45.7\%$ & $96.2\%$ / $18.68$ \\
speed-v4 & $18.57$ & $15.2\%$ & $98.0\%$ / $18.80$ \\
\bottomrule
\end{tabular}
\end{table}

The lesson is not that calibration is wrong but that an evaluation serves two
purposes---estimating performance and \emph{discriminating between
candidates}---and calibrating away the diversity destroys the second. The
perturbations a faithful single-vehicle profile removes are exactly the ones
under which a fragile policy separates from a robust one. We now require
agreement between the reset-diverse screen, targeted stress checks and physical
logs before promoting a policy. The honest limit: that screen did not predict
the speed-v2 failure either---it flagged lower success and much larger
variance, enough to rank correctly, but the failure mode was gate clearance,
which no fixed-parameter profile can surface.

\section{Results}

\textbf{Simulation under deployment-matched conditions.} Under the
reset-diverse screen the deployed policy family reached $98.0$--$98.2\%$
three-lap success at $18.32$--$18.34$\,s ($\sigma = 0.41$--$0.48$\,s). Under
the deployment-matched ground-only protocol ($1000$ ground starts) the selected
checkpoint reached $100\%$ at $18.10$\,s with $\sigma = 0.08$\,s and worst
case $18.36$\,s; two rejected alternatives also reached $100\%$ but at
$18.26$\,s and $18.59$\,s. The variance collapse between protocols is itself a
diagnostic: nearly all reset-diverse variance comes from start-state diversity
rather than from dynamics.

\begin{table}[t]
\centering
\footnotesize
\caption{Reward variants targeting lap time, under the reset-diverse screen
($768$ envs for base, $1000$ otherwise). No variant improved mean time at
matched success rate.}
\label{tab:speed}
\begin{tabular}{@{}l r r r r@{}}
\toprule
Candidate & SR (\%) & Mean (s) & $\sigma$ (s) & Best (s) \\
\midrule
base     & $98.0$ & $\mathbf{18.34}$ & $0.48$ & $17.50$ \\
speed-v1 & $86.7$ & $20.39$ & $2.12$ & $17.56$ \\
speed-v2 & $96.2$ & $18.68$ & $1.21$ & $\mathbf{17.44}$ \\
speed-v3 & $\mathbf{98.8}$ & $19.05$ & $\mathbf{0.42}$ & $18.16$ \\
speed-v4 & $98.0$ & $18.80$ & $0.43$ & $18.04$ \\
speed-v5 & $94.2$ & $21.73$ & $2.55$ & $17.90$ \\
\bottomrule
\end{tabular}
\end{table}

\textbf{Speed candidates.} Table~\ref{tab:speed} summarizes five attempts to
recover speed by reward modification under fixed dynamics and reset settings:
broad time and control pressure (v1), a post-loop velocity term on gates
$[4,5,6,0]$ (v2), the same narrowed to $[5,6]$ (v3), a velocity term without
the next-gate blend (v4), and mild time pressure alone (v5). None improved mean
time at matched success rate. The candidate with the fastest single rollout
($17.44$\,s) had the second-worst variance and failed in the arena---the
empirical basis for the selection rule above---and reintroducing a velocity
term narrowly and weakly (v3, v4) produced stable but slower policies,
suggesting the dense-to-sparse decision is not finely tunable here.

\textbf{Physical evaluation.} The deployed policy completed the timed three-lap
evaluation in $\mathbf{19.004}$\,s and continued to five correctly ordered laps
at a mean lap time of $6.236$\,s, with timing recovered from flight logs under
a strict gate-ordering rule that rejects laps containing skipped or misordered
gates. Against the deployment-matched prediction of $18.10$\,s the residual gap
is ${\sim}0.9$\,s, or $5\%$. We do not have an explanation, and want to be
precise about why: gate clearance margins, gate geometry calibration and
rate-loop overshoot beyond our first-order motor model are all consistent with
the sensitivity analysis, a fixed-parameter evaluation profile cannot
discriminate among them, and we lacked the open-loop excitation data to identify
the rate loop. The gap is open. The two alternative policies flown in the same
session ranked in the order predicted by the reset-diverse screen and against
the order predicted by the calibrated profile, which is the evidence in
Table~\ref{tab:backtest}.

\section{Discussion and Limitations}

The components most responsible for outcomes were not the ones the literature
foregrounds. The learning algorithm was standard and never a bottleneck;
network capacity was set by an interface constraint rather than a study; and the
reward mattered mostly through a few near-discontinuous thresholds---a sparse
bonus that must exceed an integrated cost, a blending coefficient with a
stability boundary---rather than through smooth tuning. What determined outcomes
were the three objects this report is organized around: the reward predicate's
relationship to the progress predicate, the initial-state distribution, and the
evaluation harness.

On randomization our experience supports a narrower prescription than the
standard one: identify the nominal first, verify it with a single rollout at the
hypothesized value before committing training time, then randomize narrowly and
spend the budget where a sensitivity analysis says the policy is fragile. For
us that was the observation map, not the transition model---the opposite of
where dynamics-randomization practice \cite{tobin2017,peng2018} directs
attention, and, we suspect, a consequence of motion capture, which removes the
perception gap and leaves a small, systematic, easily overlooked pipeline bias
in its place \cite{kaufmann2023}.

The main limitation is identification. Our estimates rest on closed-loop race
logs, which are policy-biased and not persistently exciting: they constrain
effective thrust authority and command latency adequately and say little about
motor time constant, drag, or rate-loop bandwidth. The remedy is a dedicated
open-loop identification flight with a phase-structured excitation
schedule---hover phases for mass and observation bias, vertical chirps for the
thrust curve and motor lag, per-axis rate chirps for bandwidth, straight-line
segments for drag, gate flybys for corner calibration---which we designed but
did not fly. Relatedly, we never measured the real observation bias the
sensitivity analysis identified as dominant, so those ranges are hypotheses,
and the most direct next experiment is to measure stationary velocity bias and
gate-corner error, randomize over the measured uncertainty, and re-run the
sweep with elimination of the $2.9\%$ collapse as a go/no-go criterion. We
deliberately did not pursue residual dynamics learning \cite{kaufmann2023}:
with the nominal misidentified by a factor approaching two, a residual model
would have been asked to absorb a gross parametric error, harder to fit and far
harder to debug than correcting the parameter. Finally, all results are for a
single track, vehicle and arena, and we have no evidence about how the ordering
between observation and dynamics sensitivity would change under a
perception-based state estimator, where the observation map is far noisier and
far less systematic than motion capture.

{\footnotesize
\begin{thebibliography}{99}
\itemsep 0pt

\bibitem{amodei2016} D.~Amodei et al. Concrete problems in AI safety.
\emph{arXiv:1606.06565}, 2016.

\bibitem{ferede2025} R.~Ferede et al. Simulation-to-reality gap in
reinforcement learning for agile flight: robustness versus specialization.
\emph{arXiv:2504.21586}, 2025.

\bibitem{kaufmann2023} E.~Kaufmann et al. Champion-level drone racing using
deep reinforcement learning. \emph{Nature}, 620:982--987, 2023.

\bibitem{loquercio2020} A.~Loquercio et al. Deep drone racing: from simulation
to reality with domain randomization. \emph{IEEE Trans.\ Robotics},
36(1):1--14, 2020.

\bibitem{mittal2023} M.~Mittal et al. Orbit: a unified simulation framework
for interactive robot learning environments. \emph{IEEE RA-L},
8(6):3740--3747, 2023.

\bibitem{ng1999} A.~Y.~Ng, D.~Harada, and S.~Russell. Policy invariance under
reward transformations: theory and application to reward shaping. In
\emph{ICML}, 1999.

\bibitem{pasumarti2025} P.~Pasumarti et al. Sparse-reward reinforcement
learning for autonomous drone racing. \emph{arXiv:2512.11781}, 2025.

\bibitem{peng2018} X.~B.~Peng et al. Sim-to-real transfer of robotic control
with dynamics randomization. In \emph{ICRA}, 2018.

\bibitem{rudin2022} N.~Rudin et al. Learning to walk in minutes using
massively parallel deep reinforcement learning. In \emph{CoRL}, 2022.

\bibitem{schulman2016} J.~Schulman et al. High-dimensional continuous control
using generalized advantage estimation. In \emph{ICLR}, 2016.

\bibitem{schulman2017} J.~Schulman et al. Proximal policy optimization
algorithms. \emph{arXiv:1707.06347}, 2017.

\bibitem{song2023} Y.~Song et al. Reaching the limit in autonomous racing:
optimal control versus reinforcement learning. \emph{Science Robotics}, 8(82),
2023.

\bibitem{tobin2017} J.~Tobin et al. Domain randomization for transferring deep
neural networks from simulation to the real world. In \emph{IROS}, 2017.

\end{thebibliography}
}

\end{document}
